% =========================================================
% Open Formulas, Definable Relations, and Substitution Lemmas
% =========================================================

\subsection{Semantics: Open Formulas, Definable Relations, and Substitution Lemmas}

\begin{remark*}[Section toolkit: Predicate vocabulary]
This section separates three uses that are often blurred. A predicate symbol
is object-language syntax from the signature. An open formula is a formula
with free variables. A definable relation is the set-theoretic relation in a
structure determined by an open formula.
\end{remark*}

% The legacy label def:predicate is preserved on the closest equivalent item:
% the syntactic open-formula notion previously called a predicate.
\begin{definitionbox}{Definition (Open Formula)}
\begin{definition}[Open Formula]
\label{def:predicate}
An \emph{open formula} is a well-formed formula with at least one free
variable. When the free variables are among \(x_1,\ldots,x_n\), it may be
displayed as \(\varphi(x_1,\ldots,x_n)\).
\end{definition}
\end{definitionbox}

\begin{remark*}[Standard quantified statement]
\[
\varphi \text{ is open}
\Longleftrightarrow
\varphi\in\mathsf{Form}_{\mathcal L}
\text{ and }
\mathsf{FV}(\varphi)\neq\varnothing.
\]
\end{remark*}

\begin{remark*}[Predicate reading]
\(\operatorname{OpenFormula}(\varphi)\) means that \(\varphi\) is a
well-formed formula with at least one free variable.
\end{remark*}

\begin{remark*}[Negated quantified statement]
\[
\varphi\notin\mathsf{Form}_{\mathcal L}
\quad\text{or}\quad
\mathsf{FV}(\varphi)=\varnothing.
\]
\end{remark*}

\begin{remark*}[Negation predicate reading]
\(\neg\operatorname{OpenFormula}(\varphi)\) means that \(\varphi\) is either
not well formed or is a sentence.
\end{remark*}

\begin{remark*}[Failure modes]
\begin{description}
  \item[Exposition.] An expression fails to be an open formula when it is not a well-formed formula
or when all of its variables are bound.
\end{description}
\end{remark*}

\begin{remark*}[Failure modes]
\begin{description}
  \item[Exposition.] \[
\underbrace{\varphi\notin\mathsf{Form}_{\mathcal L}}_{\text{not well formed}}
\quad\lor\quad
\underbrace{\mathsf{FV}(\varphi)=\varnothing}_{\text{sentence, not open}}.
\]
\end{description}
\end{remark*}

\begin{remark*}[Interpretation]
Open formulas are syntactic predicates in the lightest Volume I sense. They do
not by themselves name a set-theoretic relation until a structure is fixed.
\end{remark*}

\begin{remark*}[Examples]
\(P(x)\), \(x<y\), and \(\exists z\,R(x,z)\) are open formulas when the
displayed variables remain free as indicated.
\end{remark*}

\begin{remark*}[Non-Examples]
\(\forall x\,P(x)\) is a sentence, not an open formula. A predicate symbol
\(P\) alone is a symbol, not a formula.
\end{remark*}

\begin{dependencies}
\begin{itemize}
  \item \hyperref[def:predicate-symbol]{Predicate Symbol}
  \item \hyperref[def:wff-pred]{Well-Formed Formula}
\end{itemize}
\end{dependencies}

\begin{definitionbox}{Definition (Definable Relation)}
\begin{definition}[Definable Relation]
\label{def:definable-relation}
Let \(\mathcal M=\langle D,I\rangle\) be a structure and let
\(\varphi(x_1,\ldots,x_n)\) be an open formula. The \emph{relation defined by}
\(\varphi\) in \(\mathcal M\) is the set
\[
R_{\varphi}^{\mathcal M}
=
\{(a_1,\ldots,a_n)\in D^n:
\mathcal M,s[x_1\mapsto a_1,\ldots,x_n\mapsto a_n]\models\varphi\}.
\]
\end{definition}
\end{definitionbox}

\begin{remark*}[Standard quantified statement]
\[
(a_1,\ldots,a_n)\in R_{\varphi}^{\mathcal M}
\Longleftrightarrow
\mathcal M,s[x_1\mapsto a_1,\ldots,x_n\mapsto a_n]\models\varphi.
\]
\end{remark*}

\begin{remark*}[Predicate reading]
\(\operatorname{DefinableRelation}(R,\varphi,M)\) means that \(R\) is the
relation on the domain of \(M\) determined by \(\varphi\) in \(M\).
\end{remark*}

\begin{remark*}[Negated quantified statement]
\[
R\neq R_{\varphi}^{\mathcal M}.
\]
\end{remark*}

\begin{remark*}[Negation predicate reading]
\(\neg\operatorname{DefinableRelation}(R,\varphi,M)\) means that \(R\) is not
the relation determined by \(\varphi\) in \(M\).
\end{remark*}

\begin{remark*}[Failure modes]
\begin{description}
  \item[Exposition.] A proposed definable relation fails when at least one tuple is included or
excluded differently from the satisfaction condition imposed by the formula.
\end{description}
\end{remark*}

\begin{remark*}[Failure modes]
\begin{description}
  \item[Exposition.] \[
\underbrace{(a_1,\ldots,a_n)\in R \land
\mathcal M,s[\vec x\mapsto\vec a]\not\models\varphi}_{\text{extra tuple}}
\quad\lor\quad
\underbrace{(a_1,\ldots,a_n)\notin R \land
\mathcal M,s[\vec x\mapsto\vec a]\models\varphi}_{\text{missing tuple}}.
\]
\end{description}
\end{remark*}

\begin{remark*}[Interpretation]
A relation is a set-theoretic subset of \(D^n\). A definable relation is one
obtained from a formula after a structure supplies semantic meaning.
\end{remark*}

\begin{dependencies}
\begin{itemize}
  \item \hyperref[def:predicate]{Open Formula}
  \item \hyperref[def:structure]{Structure}
  \item \hyperref[def:fol-satisfaction]{Satisfaction}
\end{itemize}
\end{dependencies}

\begin{lemmabox}{Lemma (Substitution Lemma for Terms)}
\begin{lemma}[Substitution Lemma for Terms]
\label{lem:subst-terms}
Let \(\mathcal M\) be a structure, let \(s\) be a variable assignment, let
\(x\) be a variable, and let \(t\) and \(u\) be terms. Then
\[
\llbracket t[u/x]\rrbracket_{\mathcal M,s}
=
\llbracket t\rrbracket_{\mathcal M,\,
s[x\mapsto\llbracket u\rrbracket_{\mathcal M,s}]}.
\]
\hyperref[prf:subst-terms]{Go to proof.}
\end{lemma}
\end{lemmabox}

\begin{remark*}[Standard quantified statement]
\[
\forall t\,\forall u\quad
\llbracket t[u/x]\rrbracket_{\mathcal M,s}
=
\llbracket t\rrbracket_{\mathcal M,\,
s[x\mapsto\llbracket u\rrbracket_{\mathcal M,s}]}.
\]
\end{remark*}

\begin{remark*}[Predicate reading]
Term substitution commutes with semantic term evaluation in a fixed structure
and assignment.
\end{remark*}

\begin{remark*}[Interpretation]
The lemma says that syntactically substituting a term before evaluation gives
the same domain element as first evaluating the substituting term and then
updating the assignment.
\end{remark*}

\begin{dependencies}
\begin{itemize}
  \item \hyperref[def:subst-notation]{Substitution Notation}
  \item \hyperref[def:free-for]{Free for Substitution}
  \item \hyperref[def:term-interp]{Term Evaluation}
\end{itemize}
\end{dependencies}

\begin{lemmabox}{Lemma (Substitution Lemma for Formulas)}
\begin{lemma}[Substitution Lemma for Formulas]
\label{lem:subst-formulas}
Let \(\varphi\) be a formula and let \(u\) be a term free for \(x\) in
\(\varphi\). For every structure \(\mathcal M\) and assignment \(s\),
\[
\mathcal M,s\models\varphi[u/x]
\Longleftrightarrow
\mathcal M,s[x\mapsto\llbracket u\rrbracket_{\mathcal M,s}]\models\varphi.
\]
\hyperref[prf:subst-formulas]{Go to proof.}
\end{lemma}
\end{lemmabox}

\begin{remark*}[Standard quantified statement]
\[
\forall \varphi\,\forall u\quad
u \text{ free for } x \text{ in } \varphi
\Rightarrow
\bigl(
\mathcal M,s\models\varphi[u/x]
\Longleftrightarrow
\mathcal M,s[x\mapsto\llbracket u\rrbracket_{\mathcal M,s}]\models\varphi
\bigr).
\]
\end{remark*}

\begin{remark*}[Predicate reading]
Formula substitution commutes with satisfaction when the substituted term is
free for the variable being replaced.
\end{remark*}

\begin{remark*}[Interpretation]
This lemma connects the syntactic substitution operation from the syntax
section with the semantic modified-assignment operation from the semantics
section. The free-for condition prevents variable capture.
\end{remark*}

\begin{dependencies}
\begin{itemize}
  \item \hyperref[def:subst-notation]{Substitution Notation}
  \item \hyperref[def:free-for]{Free for Substitution}
  \item \hyperref[def:fol-satisfaction]{Satisfaction}
  \item \hyperref[lem:subst-terms]{Substitution Lemma for Terms}
\end{itemize}
\end{dependencies}
